\subsection{商模的对偶。直和的对偶。对偶基}

我们将第1号的定理1应用于 \(\mathbf{F} = \mathbf{A}_{\mathrm{s}}\) 的情形：

命题9。设 \(\mathbf{E}'\) , \(\mathbf{E}\) , \(\mathbf{E}''\) 是A-模，且

\[
\mathrm{E} ^ {\prime} \stackrel {u} {\longrightarrow} \mathrm{E} \stackrel {v} {\longrightarrow} \mathrm{E} ^ {\prime \prime} \longrightarrow 0\tag{21}
\]

线性映射的一个正合序列。则转置映射的序列

\[
0 \longrightarrow \mathrm{E} ^ {\prime \prime *} \stackrel {t _ {v}} {\longrightarrow} \mathrm{E} ^ {*} \stackrel {t _ {u}} {\longrightarrow} \mathrm{E} ^ {\prime *} \tag {1}
\]

是正合的。

推论。设M是A-模E的一个子模，且 \(\phi: \mathbf{E} \to \mathbf{E} / \mathbf{M}\) 为典范同态。则 \(^t\phi\) 是E/M的对偶到E*中与M正交的子模M'上的一个同构。

若 \(j: \mathbf{M} \to \mathbf{E}\) 是典范单射，则 \(j\) 的核按定义是 \(\mathbf{M}\) 在 \(\mathbf{E}^*\) 中的正交。

此外，在推论的记号下，一个单同态 \(E^{*}/M^{\prime} \rightarrow M^{*}\) 可由 \(^{t}j\) 过渡到商而得到。

命题10. 设 \((\mathbf{E}_t)_{t \in I}\) 是一族A-模，且对所有 \(t \in I\) ，设 \(j_t: \mathbf{E}_t \to \mathbf{E} = \bigoplus_{t \in I} \mathbf{E}_t\) 为典范单射。则乘积映射 \(x^* \mapsto (t j_t(x^*))\) 是 \(\mathbf{E}^*\) （ \(\mathbf{E}\) 的对偶）到乘积 \(\prod_{t \in I} \mathbf{E}_t^*\) 上的一个同构。

这是§1第6号，命题6的推论1的一个特例，应用于以下情形： \(\prod_{\lambda} F_{\lambda} = A_{s}\) .

如果借助典范单射 \(j_{i}\) 把诸 \(E_{i}\) 与它们的直和 \(E\) 的子模等同起来，并且借助乘积映射 \(x^{*} \mapsto (t j_{i}(x^{*}))\) 把 \(E^{*}\) 与 \(\prod_{i \in I} E_{i}^{*}\) 等同起来，那么可以说 \(\prod_{i \in I} E_{i}^{*}\) 是 \(\bigoplus_{i \in I} E_{i}\) 的对偶，其中典范双线性形式由

\[
\langle (x _ {i}), (x _ {i} ^ {*}) \rangle = \sum_ {i \in 1} \langle x _ {i}, x _ {i} ^ {*} \rangle .\tag{22}
\]

推论. 设 M, N 是A-模 E 中两个互补的子模，且 \(p:E\to M\) , \(q:E\to N\) 是相应的投影算子；则 \({}^{tp}+{}^{tq}:M^{*}\oplus N^{*}\to E^{*}\) 是一个同构，且 \({}^{tp}\) （相应地 \({}^{tq}\) ）是 \(M^{*}\) （相应地 \(N^{*}\) ）到 \(E^{*}\) 中与 N（相应地 M）正交的子模上的一个同构。此外，若 \(i:M\to E\) 和 \(j:N\to E\) 是典范单射，则 \({}^{tp}\circ{}^{ti}\) 和 \({}^{tq}\circ{}^{tj}\) 是投影算子 \(E^{*}\to{}^{tp}(M^{*})\) , \(E^{*}\to{}^{tq}(N^{*})\) ，它们对应于把 \(E^{*}\) 分解为直和 \({}^{tp}(M^{*})\) 与 \({}^{tq}(N^{*})\) 。

\(p \circ i = 1_{\mathbf{M}}, \quad q \circ j = 1_{\mathbf{N}}, \quad p \circ j = q \circ i = 0, \quad i \circ p + j \circ q = 1_{\mathbf{E}},\) 由此，通过转置（第5号，公式（16）、（17）和（18））， \(^{ti} \circ ^{tp} = 1_{\mathbf{M}^*}, \; ^{tj} \circ ^{tq} = 1_{\mathbf{N}^*}, \; ^{tj} \circ ^{tp} = ^{ti} \circ ^{tq} = 0. \; ^{tp} \circ ^{ti} + ^{tq} \circ ^{tj} = 1_{\mathbf{E}^*}\) 且该命题由§1第6号命题6的推论2得出。

在推论的假设下， \(M^{*}\) （相应地 \(N^{*}\) ）通常与正交 \(^{tp}(M^{*})\) （相应地 \(^{t}q(N^{*})\) ）——即N（相应地M）在 \(E^{*}\) 中的正交——等同起来，从而把M（相应地N）上的每个线性形式u等同于E上扩张u且在N（相应地M）上为零的那个线性形式。

当A-模E具有一个基 \((e_{t})_{t\in\mathbb{T}}\) 时，已经看到，给出这个基典范地定义了一个同构 \(u:A_{s}^{\mathrm{(T)}}\to E\) 。根据命题10和第3节命题5， \(\mathbf{A}_{s}^{\mathrm{(T)}}\) 与乘积典范等同 \(A_{d}^{T}\) 的对偶；考虑逆步同构 \(\check{u}:A_{d}^{T}\to E^{*}\) 的典范映射。如果对所有 \(t\in T,f_{t}\) 是……的元素 \(A_{d}^{T}\) 其所有投影均为零，除了指标为t的那个等于1，并且如果我们记 \(e_{t}^{*}=\check{u}(f_{t})\) ，元素 \(e_{t}^{*}\) 于 \(E^{*}\) 则由（19）和（22）刻画的关系为

\[
\langle e _ {t}, e _ {t ^ {\prime}} ^ {*} \rangle = \left\{ \begin{array}{l l} 0 & \text { for } t ^ {\prime} \neq t \\ 1 & \text { for } t ^ {\prime} = t. \end{array} \right.\tag{23}
\]

这等同于说，对于所有 \(x = \sum_{t \in T} \xi_t e_t \in E\) , \(e_t^*(x) = \xi_t\) ；也称 \(e_t^*\) 为指标 \(t\) 在 \(E\) 上的坐标形式。由(23)可知， \((e_t^*)\) 是 \(E^*\) 中的一个自由系。

特别地，若T是有限的，则诸 \(e_{t}^{*}\) 构成 \(E^{*}\) 的一组基，诸 \(f_{t}\) 此时构成 \(A_{d}^{T}\) 的典范基。因此：

命题11. 具有由 \(n\) 个元素组成的一组基的自由模，其对偶是具有由 \(n\) 个元素组成的一组基的自由模。

注意，具有无限基的自由模的对偶模不一定是自由模（VII，§3，习题10）。

定义7. 若E是具有有限基 \((e_{t})\) 的自由模，则基 \((e_{t}^{*})\) ——即E的对偶 \(E^{*}\) 的、由关系式(23)所定义的基——称为 \((e_{t})\) 的对偶基。

关系式(23)也可以写成如下形式

\[
\langle e _ {t}, e _ {t ^ {\prime}} ^ {*} \rangle = \delta_ {t t ^ {\prime}}\tag{24}
\]

其中 \(\delta_{tt'}\) 是 \(\mathbf{T} \times \mathbf{T}\) 上的克罗内克符号。

注意，如果T是有限的且 \((e_{t}^{*})\) 是 \((e_{t})\) 的对偶基，则对于 \(x = \sum_{t \in T} \xi_{t} e_{t} \in E\) , \(x^{*} = \sum_{t \in T} \xi_{t}^{*} e_{t}^{*} \in E^{*}\) ,

\[
\langle x, x ^ {*} \rangle = \sum_ {t \in T} \xi_ {t} \xi_ {t} ^ {*}.\tag{25}
\]

右A-模的有限基的对偶基当然也类似地定义。

推论。有限生成投射模的对偶是有限生成投射模。

有限生成的投射左A-模可以等同于具有有限基的自由A-模 \(A_{s}^{n}\) 的一个直因子M（第2号，命题4的推论1）。于是（命题11及命题10的推论） \(M^{*}\) 同构于 \(A_{d}^{n}\) 的一个直因子，由此得出推论。

命题12. 设 E 是一个 A-模，且 \((a_{t})_{t\in \mathbf{T}}\) 是 E 的一个生成系。以下条件等价：

\begin{enumerate}
\def\labelenumi{(\alph{enumi})}
\item
  E 是投射 A-模。
\item
  存在一个族 \((a_{t}^{*})_{t\in \mathbf{T}}\) （由 \(\mathbf{E}\) 上的线性形式组成），使得对所有 \(x\in \mathbf{E}\) ，族 \((\langle x,a_t^*\rangle)_{t\in \mathbf{T}}\) 具有有限支撑，且
\end{enumerate}

\[
x = \sum_ {t \in \mathbf {T}} \langle x, a _ {t} ^ {*} \rangle a _ {t}.\tag{26}
\]

存在一个满同态 \(u: \mathbf{L} \to \mathbf{E}\) ，其中 \(\mathbf{L} = \mathbf{A}_s^{(\mathrm{T})}\) ，使得若 \((e_t)_{t \in \mathbb{T}}\) 是 \(\mathbf{L}\) 的典范基，则 \(u(e_t) = a_t\) （§ 1，第 11 号，命题 17）；为使 \(\mathbf{E}\) 是投射的，必要且充分的是存在线性映射 \(v: \mathbf{E} \to \mathbf{L}\) 使得 \(u \circ v = 1_{\mathbb{E}}\) （第2号，命题4及§1，第9号，命题15）。如果这样的映射存在，并记 \(^tv(e_t^*) = a_t^*\) ，则 \(\langle x, a_t^* \rangle = \langle x, ^tv(e_t^*) \rangle = \langle v(x), e_t^* \rangle\) ，从而族 \((\langle x, a_t^* \rangle)\) 具有有限支撑，且 \(x = u\left(\sum_{t \in \mathbf{T}} \langle (x), e_t^* \rangle e_t\right) = \sum_{t \in \mathbf{T}} \langle x, a_t^* \rangle a_t\) 对所有 \(x \in \mathbf{E}\) 成立。反之，若命题的条件（b）成立，则和 \(\sum_{t \in \mathbf{T}} \langle x, a_t^* \rangle e_t\) 对所有 \(x \in \mathbf{E}\) 有定义，且 \(x \to \sum_{t \in \mathbf{T}} \langle x, a_t^* \rangle e_t\) 是一个线性映射 \(v: \mathbf{E} \to \mathbf{L}\) ，使得 \(u \circ v = 1_{\mathbf{E}}\) 。

